\chapter{Wzór Itô: rachunek różniczkowy z wariacją kwadratową}
\label{ch:ito-formula}

\section{Motywacja: dlaczego zwykła reguła łańcuchowa zawodzi?}
Niech $W$ będzie procesem Wienera względem filtracji ustalonej
we wprowadzeniu. Gdybyśmy mogli bez zmian stosować klasyczną
regułę łańcuchową, uzyskalibyśmy
$W_t^2=2\int_0^tW_s\,\dd W_s$. Tymczasem w rozdziale
\ref{ch:ito-construction} wyprowadziliśmy ścisłą tożsamość
$W_t^2=2\int_0^tW_s\,\dd W_s+t$.
Dodatkowy składnik nie jest błędem przybliżenia: wynika z tego,
że suma kwadratów przyrostów $W$ nie znika w granicy, lecz
zbiega do długości przedziału.

W rachunku deterministycznym, przy różniczkowaniu $F(x)$,
wyraz drugiego rzędu jest rzędu $(\Delta x)^2$ i po
zsumowaniu zanika dla trajektorii o skończonej wariacji.
Dla ruchu Browna typowy przyrost na odcinku długości
$\Delta t$ jest rzędu $\sqrt{\Delta t}$, a więc jego
kwadrat jest rzędu $\Delta t$. Suma takich wyrazów nie
zanika. Celem rozdziału jest zastąpienie tej heurystyki
ścisłym twierdzeniem.

\section{Proces Itô i jego wariacja kwadratowa}
\paragraph{Konwencja zbieżności.}
Dla procesów $Y^{(n)}$ i $Y$ na $[0,T]$ piszemy
$Y^{(n)}\to Y$ \emph{ucp} (jednostajnie na przedziałach
zwartych według prawdopodobieństwa), jeżeli dla każdego
$\varepsilon>0$ zachodzi
$\Prob(\sup_{0\le t\le T}|Y_t^{(n)}-Y_t|>\varepsilon)
\to0$; gdy czas przebiega $[0,\infty)$, wymaga się
tej własności dla każdego skończonego $T$.

\begin{definicja}[Proces Itô]
Niech $T>0$, $W$ będzie rzeczywistym procesem Wienera,
$X_0$ będzie rzeczywistą zmienną $\calF_0$-mierzalną,
$b$ będzie procesem progresywnie mierzalnym, a $\sigma$
procesem przewidywalnym. Zakładamy, że prawie na pewno
$\int_0^T|b_s|\,\dd s<\infty$ oraz
$\int_0^T\sigma_s^2\,\dd s<\infty$.
Proces
\[
X_t=X_0+\int_0^tb_s\,\dd s+\int_0^t\sigma_s\,\dd W_s,
\qquad 0\le t\le T,
\tag{3.1}\label{eq:ito-process}
\]
nazywamy \emph{procesem Itô}. Druga całka jest całką
skonstruowaną w rozdziale \ref{ch:ito-construction}
(przez lokalizację, jeśli to konieczne).
\end{definicja}

\begin{uwaga}[Dlaczego nie zakładamy skończonych wartości oczekiwanych?]
Do dowodu lokalnego wzoru Itô wystarcza skończoność całek
z poprzedniej definicji dla niemal każdej trajektorii.
Zatrzymując proces, możemy na każdym ograniczonym przedziale
kontrolować zarówno $|X|$, jak i energię $\int\sigma^2$,
zmienność $\int|b|$ oraz pochodne funkcji, do której
stosujemy wzór. Dopiero po wykonaniu rachunku zdejmujemy
zatrzymanie. Nie wolno zatem bez uzasadnienia zakładać, że
całki występujące we wzorze mają skończone wartości oczekiwane.
\end{uwaga}

\begin{lemat}[Wariacja kwadratowa procesu Itô]
\label{lem:ito-qv}
Niech $X$ będzie procesem \eqref{eq:ito-process}. Dla
każdego ciągu deterministycznych podziałów $\pi_n$ przedziału
$[0,T]$ o średnicy dążącej do zera sumy
$\sum_{[u,v]\in\pi_n,\,v\le t}(X_v-X_u)^2$
zbiegają do $\int_0^t\sigma_s^2\,\dd s$ jednostajnie
w prawdopodobieństwie na $[0,T]$ (w skrócie: ucp).
W szczególności $[X]_t=\int_0^t\sigma_s^2\,\dd s$.
\end{lemat}
\begin{proof}
Piszemy $X=X_0+A+M$, gdzie
$A_t=\int_0^tb_s\,\dd s$ jest ciągłym procesem
skończonej wariacji, a $M_t=\int_0^t\sigma_s\,\dd W_s$
ciągłym lokalnym martyngałem. Na dowolnym podziale mamy
\[
\sum(\Delta X)^2-\sum(\Delta M)^2
=2\sum\Delta A\,\Delta M+\sum(\Delta A)^2.
\]
Dla ustalonej trajektorii ciągłość $A$ zapewnia
$\max|\Delta A|\to0$, a stąd
\[
\sum|\Delta A|^2
\le (\max|\Delta A|)\operatorname{TV}_{[0,T]}(A)
\longrightarrow0.
\]
Z nierówności Cauchy'ego--Schwarza otrzymujemy
\[
\left|\sum\Delta A\,\Delta M\right|
\le\Big(\sum(\Delta A)^2\Big)^{1/2}
\Big(\sum(\Delta M)^2\Big)^{1/2}.
\]
Drugi czynnik jest ograniczony w prawdopodobieństwie,
ponieważ sumy kwadratów przyrostów ciągłego lokalnego
martyngału zbiegają ucp do $[M]_T$; wynik ten udowodniono
w tomie II. Oszacowania stosują się również do sum
uciętych w dowolnej chwili $t$, z supremum po $t$.
Otrzymujemy więc zbieżność ucp sum dla $X$ i $M$ do
tej samej granicy. Z rozdziału \ref{ch:ito-construction}
wiemy, że $[M]_t=\int_0^t\sigma_s^2\,\dd s$.
\end{proof}

\begin{lemat}[Ważone sumy kwadratów przyrostów]
\label{lem:weighted-qv}
Niech $X$ będzie ciągłym procesem o wariacji kwadratowej
$[X]$ ciągłej i niemalejącej, przy czym sumy kwadratów
przyrostów na deterministycznych podziałach o średnicy
dążącej do zera zbiegają do $[X]$ ucp. Niech $G$ będzie
ciągłym procesem o trajektoriach ograniczonych na $[0,T]$.
Wówczas dla takich podziałów
\[
\sum_{[u,v]\in\pi_n,\,v\le t}G_u(X_v-X_u)^2
\ \xrightarrow[\,n\to\infty\,]{\mathrm{ucp}}\
\int_0^tG_s\,\dd[X]_s.
\tag{3.2}\label{eq:weighted-qv}
\]
Po prawej stronie jest zwykła całka Lebesgue'a--Stieltjesa
względem niemalejącej, ciągłej trajektorii $[X]$.
\end{lemat}
\begin{proof}
Ustalmy najpierw deterministyczny podział pomocniczy
$0=r_0<\cdots<r_m=T$ i zastąpmy $G$ procesem
$G^{(m)}_s=G_{r_j}$ dla $s\in(r_j,r_{j+1}]$.
Dla takiej funkcji wynik wynika z ucp-zbieżności
sum kwadratów na kolejnych podprzedziałach i z
ciągłości $[X]$. Jeżeli siatka $\pi_n$ nie zawiera
wszystkich $r_j$, dołączamy je: wpływ skończenie wielu
rozdzielonych przyrostów jest zaniedbywalny, gdyż
$\sup|\Delta X|\to0$, a całkowite sumy kwadratów są
ograniczone w prawdopodobieństwie.

Pozostaje oszacować błąd zastąpienia $G$ przez $G^{(m)}$.
Jego moduł w supremum po $t$ nie przekracza
$\norm{G-G^{(m)}}_\infty\sum_{[u,v]\in\pi_n}(X_v-X_u)^2$,
a dla całki granicznej jest ograniczony przez
$\norm{G-G^{(m)}}_\infty[X]_T$. Z ciągłości trajektorii
$G$ otrzymujemy $\norm{G-G^{(m)}}_\infty\to0$ prawie
na pewno przy zagęszczaniu pomocniczej siatki. Ponadto
$[X]_T$ i sumy kwadratów są ograniczone w prawdopodobieństwie.
Najpierw wybieramy $m$ dostatecznie duże, a następnie $n$;
daje to \eqref{eq:weighted-qv}.
\end{proof}

\section{Jednowymiarowy wzór Itô}
\begin{twierdzenie}[Wzór Itô dla funkcji klasy $C^{1,2}$]
\label{th:ito1}
Niech $X$ będzie procesem \eqref{eq:ito-process}, a
$F:[0,T]\times\R\to\R$ funkcją klasy $C^{1,2}$,
tzn. $F$, $\partial_tF$, $\partial_xF$ i
$\partial_{xx}F$ są ciągłe na $[0,T]\times\R$
(jednostronnie na końcach przedziału czasu). Wtedy
z prawdopodobieństwem jeden, jednocześnie dla wszystkich
$t\in[0,T]$, zachodzi
\[
\begin{split}
F(t,X_t)=F(0,X_0)
&+\int_0^t\partial_tF(s,X_s)\,\dd s
+\int_0^t\partial_xF(s,X_s)b_s\,\dd s\\
&+\int_0^t\partial_xF(s,X_s)\sigma_s\,\dd W_s
+\frac12\int_0^t\partial_{xx}F(s,X_s)\sigma_s^2\,\dd s.
\end{split}
\tag{3.3}\label{eq:ito1}
\]
Stochastyczna całka po prawej stronie jest rozumiana lokalnie.
\end{twierdzenie}
\begin{proof}
\emph{Krok 1: lokalizacja.} Dla liczb całkowitych $m\ge1$
zatrzymajmy proces przy wyjściu $X$ z $[-m,m]$ lub
przekroczeniu przez $\int_0^t(|b_s|+\sigma_s^2)\,\dd s$
poziomu $m$. Ciągłość trajektorii i skończoność całek
zapewniają, że czasy zatrzymania rosną do $T$ prawie na pewno.
Na zbiorze zatrzymanych trajektorii $X$ pozostaje w
przedziale zwartym, a $\partial_tF$, $\partial_xF$,
$\partial_{xx}F$ są tam ograniczone i jednostajnie ciągłe.
Wystarczy wykazać tożsamość do każdego z czasów zatrzymania
i następnie usunąć lokalizację. Możemy więc w dalszej
części dowodu pracować z ograniczonym zakresem $X$
i ze skończonymi, kontrolowanymi energiami.

\emph{Krok 2: rozwinięcie na podziale.} Weźmy podział
$\pi_n=\{0=t_0<\dots<t_{N_n}=t\}$ o średnicy
dążącej do zera. Oznaczmy $\Delta_kX=X_{t_{k+1}}-X_{t_k}$
i $\Delta t_k=t_{k+1}-t_k$. Rozbijamy przyrost funkcji na zmianę wartości przestrzennej
w chwili $t_k$ oraz zmianę czasu przy 
ustalonej \emph{nowej} wartości przestrzennej:
\[
F(t_{k+1},X_{t_{k+1}})-F(t_k,X_{t_k})
=\bigl[F(t_k,X_{t_{k+1}})-F(t_k,X_{t_k})\bigr]
 +\bigl[F(t_{k+1},X_{t_{k+1}})-F(t_k,X_{t_{k+1}})\bigr].
\]
Jednostajna ciągłość pochodnych na odpowiednim zwartym
prostokącie i wzór Taylora dają
\[
\begin{split}
F(t_{k+1},X_{t_{k+1}})-F(t_k,X_{t_k})
={}&\partial_tF(t_k,X_{t_k})\Delta t_k
+\partial_xF(t_k,X_{t_k})\Delta_kX\\
&+\frac12\partial_{xx}F(t_k,X_{t_k})(\Delta_kX)^2+r_{k,n},
\end{split}
\]
gdzie suma reszt $\sum_k r_{k,n}$ dąży do zera
w prawdopodobieństwie. Uzasadnijmy to bez ukrytego
założenia o istnieniu pochodnej mieszanej $F_{tx}$.
Niech $\omega_t(\delta)$ i $\omega_{xx}(\delta)$
będą modułami ciągłości, odpowiednio, $F_t$ i $F_{xx}$
na ustalonym zwartym prostokącie zawierającym
zatrzymaną trajektorię. Z twierdzenia o wartości
średniej w zmiennej $t$ i wzoru Taylora w zmiennej $x$
wynika
\[
\sum_k|r_{k,n}|\le
T\,\omega_t\!\left(|\pi_n|+\max_k|\Delta_kX|\right)
+\tfrac12\omega_{xx}\!\left(\max_k|\Delta_kX|\right)
\sum_k(\Delta_kX)^2.
\]
Rzeczywiście, w składniku czasowym odejmujemy
$F_t(t_k,X_{t_k})\Delta t_k$ od
$\int_{t_k}^{t_{k+1}}F_t(s,X_{t_{k+1}})\,\dd s$;
różnica jest kontrolowana przez pierwszy moduł ciągłości.
W składniku przestrzennym stosujemy jednowymiarowy
wzór Taylora z resztą kontrolowaną przez drugi moduł.
Z ciągłości $X$ mamy $\max_k|\Delta_kX|\to0$
prawie na pewno, a z lematu \ref{lem:ito-qv} sumy
kwadratów przyrostów są ograniczone w prawdopodobieństwie.
Oba składniki oszacowania dążą więc do zera w
prawdopodobieństwie. To rozumowanie wymaga dokładnie
klasy $C^{1,2}$, a nie silniejszej regularności.

\emph{Krok 3: identyfikacja granic.} Sumy
$\sum_k\partial_tF(t_k,X_{t_k})\Delta t_k$
zbiegają do $\int_0^t\partial_tF(s,X_s)\,\dd s$
ścieżkowo, ponieważ integrand jest ciągły.
Dla części liniowej piszemy
$\Delta_kX=\int_{t_k}^{t_{k+1}}b_s\,\dd s
+\int_{t_k}^{t_{k+1}}\sigma_s\,\dd W_s$.
Funkcje schodkowe
$H_s^{(n)}=\partial_xF(t_k,X_{t_k})$ dla
$s\in(t_k,t_{k+1}]$ są przewidywalne. Ich różnica
od ciągłego, adaptowanego procesu
$H_s=\partial_xF(s,X_s)$ dąży jednostajnie do zera
na każdej zatrzymanej trajektorii. Zatem
$\int H^{(n)}b\,\dd s\to\int Hb\,\dd s$
ścieżkowo, a po odpowiedniej lokalizacji
$\E\int|H^{(n)}-H|^2\sigma^2\,\dd s\to0$ przez
zbieżność zdominowaną. Izometria Itô wraz z
nierównością maksymalną Dooba daje zbieżność
odpowiednich całek stochastycznych.

Dla wyrazów kwadratowych stosujemy lemat
\ref{lem:weighted-qv} z
$G_s=\partial_{xx}F(s,X_s)$, otrzymując
\[
\sum_kG_{t_k}(\Delta_kX)^2
\longrightarrow\int_0^tG_s\sigma_s^2\,\dd s
\]
w prawdopodobieństwie. Sumując wszystkie granice,
otrzymujemy \eqref{eq:ito1} dla ustalonego $t$.
Obie strony mają ciągłe trajektorie. Równość w chwilach
wymiernych (oraz w $T$) daje ją jednocześnie dla wszystkich
$t$, poza jednym zbiorem zerowym. Na końcu usuwamy
lokalizację, pozwalając poziomom zatrzymania rosnąć do nieskończoności.
\end{proof}

\begin{uwaga}[Znaczenie wzoru różniczkowego]
Często zapisujemy \eqref{eq:ito1} symbolicznie jako
$\dd F(t,X_t)=F_t\,\dd t+F_x\,\dd X_t
+\frac12F_{xx}\,\dd[X]_t$. Nie jest to równość
między zwykłymi różniczkami: jej ścisłym znaczeniem
jest tożsamość całkowa \eqref{eq:ito1}.
\end{uwaga}

\section{Przykłady: kwadrat, potęgi, wykładnik i martyngały}
\begin{zadanie}[Potęgi procesu Wienera]
Dla $n\in\N$, $n\ge2$, obliczyć rozkład na składnik
martyngałowy i składnik skończonej wariacji procesu $W_t^n$.

\emph{Rozwiązanie.} Dla $F(x)=x^n$ mamy
$F_x(x)=nx^{n-1}$ oraz $F_{xx}(x)=n(n-1)x^{n-2}$.
Wzór Itô daje
\[
W_t^n=n\int_0^tW_s^{n-1}\,\dd W_s
+\frac{n(n-1)}2\int_0^tW_s^{n-2}\,\dd s.
\]
Przy $n=2$ dostajemy $W_t^2-t=2\int_0^tW_s\,\dd W_s$;
przy $n=3$ otrzymujemy $W_t^3-3\int_0^tW_s\,\dd s
=3\int_0^tW_s^2\,\dd W_s$. Nie wolno utożsamiać
składnika całkowego z martyngałem jednostajnie
całkowalnym na całym $[0,\infty)$ bez sprawdzenia
jego momentów. Na dowolnym skończonym horyzoncie
potrzebna całkowalność wynika z momentów Gaussa.
\end{zadanie}

\begin{twierdzenie}[Wykładnik stochastyczny dla stałego parametru]
Niech $\theta\in\R$ i $W$ będzie procesem Wienera na
$[0,T]$. Wówczas $Z_t=\exp(\theta W_t-\tfrac12\theta^2t)$
jest dodatnim martyngałem i
\[
Z_t=1+\theta\int_0^t Z_s\,\dd W_s.
\tag{3.4}\label{eq:exp-mg}
\]
\end{twierdzenie}
\begin{proof}
Stosując twierdzenie \ref{th:ito1} do
$F(t,x)=\exp(\theta x-\frac12\theta^2t)$, zauważamy,
że $F_t+\frac12F_{xx}=0$, więc otrzymujemy
\eqref{eq:exp-mg} (początkowo lokalnie).
Dla $0\le s<t$ niezależny od $\calF_s$ przyrost
$W_t-W_s$ ma rozkład normalny $\Normal(0,t-s)$.
Stąd
$\E[Z_t\mid\calF_s]=Z_s\E\exp(\theta(W_t-W_s)
-\frac12\theta^2(t-s))=Z_s$.
Wszystkie wartości oczekiwane są skończone, a
$\E Z_t=1$.
\end{proof}

\begin{zadanie}[Wykładnik czasu i równanie ciepła]
Dla $u(t,x)=x^2+(T-t)$, $0\le t\le T$,
wykazać, że $u(t,W_t)$ jest martyngałem i
wyjaśnić związek z równaniem ciepła.

\emph{Rozwiązanie.} Mamy $u_t=-1$, $u_{xx}=2$,
więc $u_t+\frac12u_{xx}=0$. Wzór Itô daje
$\dd u(t,W_t)=2W_t\,\dd W_t$ i tym samym
$u(t,W_t)=T+2\int_0^tW_s\,\dd W_s$.
Jest to martyngał kwadratowo całkowalny na $[0,T]$.
W ogólności, jeśli $u\in C^{1,2}$ spełnia
$u_t+\frac12u_{xx}=0$ oraz wystarczające warunki
całkowalności, to $u(t,W_t)$ jest martyngałem.
Pokazuje to probabilistyczną rolę równania ciepła;
pełny wzór Feynmana--Kaca pojawi się później.
\end{zadanie}

\section{Iloczyny, całkowanie przez części i kowariacja}
\begin{twierdzenie}[Wzór na iloczyn procesów Itô]
\label{th:ito-product}
Niech $X,Y$ będą procesami Itô, napędzanymi tym samym
$W$, o współczynnikach odpowiednio $(b,\sigma)$ oraz
$(c,\rho)$, spełniających warunki z definicji procesu Itô.
Wówczas z prawdopodobieństwem jeden, dla wszystkich $t$,
\[
X_tY_t=X_0Y_0+\int_0^tX_sc_s\,\dd s
+\int_0^tY_sb_s\,\dd s
+\int_0^t(X_s\rho_s+Y_s\sigma_s)\,\dd W_s
+\int_0^t\sigma_s\rho_s\,\dd s.
\tag{3.5}\label{eq:ito-product}
\]
W szczególności $[X,Y]_t=\int_0^t\sigma_s\rho_s\,\dd s$.
\end{twierdzenie}
\begin{proof}
Tożsamość $4xy=(x+y)^2-(x-y)^2$ pozwala wyprowadzić
wzór na iloczyn z jednowymiarowego wzoru Itô dla
$X+Y$ i $X-Y$: stosujemy \eqref{eq:ito1} do kwadratu,
odejmujemy odpowiednie równości i dzielimy przez cztery.
Daje to składniki liniowe oraz kowariację
$\int\sigma\rho\,\dd s$. Alternatywnie, z definicji
kowariacji i polaryzacji $[X,Y]=\frac14([X+Y]-[X-Y])$,
a następnie $[X+Y]-[X-Y]=4\int\sigma\rho\,\dd s$.
\end{proof}

\begin{zadanie}[Całkowanie przez części a momenty]
Niech $W$ będzie procesem Wienera. Obliczyć
$\int_0^t sW_s\,\dd W_s$.

\emph{Rozwiązanie.} Stosujemy wzór na iloczyn do
$X_s=s$ i $Y_s=W_s^2$, dla którego
$\dd Y_s=2W_s\,\dd W_s+\dd s$.
Ponieważ $[s,W^2]_t=0$, otrzymujemy
$\dd(sW_s^2)=W_s^2\,\dd s+s(2W_s\,\dd W_s+\dd s)$.
Zatem
\[
\int_0^t sW_s\,\dd W_s
=\frac12\left(tW_t^2-\int_0^t(W_s^2+s)\,\dd s\right).
\]
Sprawdzenie: wartość oczekiwana prawej strony wynosi
$\frac12(t^2-\int_0^t2s\,\dd s)=0$, zgodnie z
martyngałowością całki.
\end{zadanie}

\section{Wzór wielowymiarowy i kierunek dyfuzji}
\begin{lemat}[Całki względem różnych współrzędnych Wienera]
\label{lem:multi-orthogonality}
Niech $W=(W^1,\ldots,W^m)$ będzie standardowym
$m$-wymiarowym procesem Wienera względem filtracji
$(\calF_t)$, to znaczy wektor jego przyrostów na
$(s,t]$ jest niezależny od $\calF_s$ i ma macierz
kowariancji $(t-s)I_m$. Jeśli $H,K$ są rzeczywistymi
procesami przewidywalnymi, dla których
$\E\int_0^T(H_s^2+K_s^2)\,\dd s<\infty$, to
\[
\E\left[
\left(\int_0^tH_s\,\dd W_s^i\right)
\left(\int_0^tK_s\,\dd W_s^j\right)\right]
=\delta_{ij}\E\int_0^tH_sK_s\,\dd s.
\]
Ponadto kowariacja ciągłych martyngałów
$M^i=\int H\,\dd W^i$ i
$N^j=\int K\,\dd W^j$ wynosi
$[M^i,N^j]_t=\delta_{ij}\int_0^tH_sK_s\,\dd s$.
W wersji lokalnie całkowalnej ostatni wzór nadal
obowiązuje dzięki lokalizacji.
\end{lemat}
\begin{proof}
Dla prostych procesów przewidywalnych przechodzimy do
wspólnego deterministycznego podziału. Iloczyn dwóch
sum zawiera wyrazy z różnych przedziałów i z tego
samego przedziału. Pierwsze mają wartość oczekiwaną
zero dzięki iterowanemu warunkowaniu; dla drugich
\[
\E[\Delta_kW^i\Delta_kW^j\mid\calF_{t_k}]
=\delta_{ij}(t_{k+1}-t_k).
\]
Jest to dokładnie podany wzór kowariancji wartości
końcowych. Przy $i\ne j$ ten sam rachunek po
warunkowaniu pokazuje, że iloczyn $M^iN^j$ jest
martyngałem. Przy $i=j$ proces
$M^iN^i-\int HK\,\dd s$ jest martyngałem,
co wynika z analogicznego rozwinięcia przyrostów.
Jednoznaczność takiej korekty wynika z faktu, że
ciągły lokalny martyngał o skończonej wariacji jest
stały: różnica dwóch korekt byłaby jednocześnie
lokalnym martyngałem i procesem skończonej wariacji.
To identyfikuje kowariację.
Dla ogólnych integrandów stosujemy gęstość procesów
prostych, izometrię Itô i nierówność Cauchy'ego--Schwarza.
W wersji lokalnej najpierw obcinamy energię obu
integrandów, a potem usuwamy zatrzymanie.
\end{proof}

\begin{twierdzenie}[Wielowymiarowy wzór Itô]
\label{th:ito-multi}
Niech $W=(W^1,\ldots,W^m)$ będzie standardowym
$m$-wymiarowym procesem Wienera względem wspólnej
filtracji. Niech $X=(X^1,\ldots,X^d)$ będzie procesem
postaci
\[
X_t^i=X_0^i+\int_0^t b_s^i\,\dd s
+\sum_{k=1}^{m}\int_0^t\sigma_s^{ik}\,\dd W_s^k,
\quad i=1,\ldots,d,
\]
gdzie $X_0$ jest $\calF_0$-mierzalny, $b$ jest
progresywnie mierzalny, $\sigma$ przewidywalna i
$\int_0^T(\norm{b_s}+\norm{\sigma_s}_{\mathrm F}^2)\,\dd s
<\infty$ prawie na pewno. Jeśli
$F\in C^{1,2}([0,T]\times\R^d)$, to z prawdopodobieństwem
jeden, dla wszystkich $t\in[0,T]$,
\[
\begin{split}
F(t,X_t)=F(0,X_0)&+\int_0^t\left(
F_t+\sum_{i=1}^d b_s^iF_{x_i}
+\frac12\sum_{i,j=1}^d(\sigma_s\sigma_s^{\mathsf T})_{ij}
F_{x_ix_j}\right)(s,X_s)\,\dd s\\
&+\sum_{k=1}^m\int_0^t
\left(\sum_{i=1}^dF_{x_i}(s,X_s)\sigma_s^{ik}\right)
\,\dd W_s^k.
\end{split}
\tag{3.6}\label{eq:ito-multi}
\]
\end{twierdzenie}
\begin{proof}
Stosujemy wielowymiarowy wzór Taylora drugiego rzędu
na deterministycznym podziale do funkcji $F(t,X_t)$.
Kontrola reszt odbywa się tak jak w dowodzie twierdzenia
\ref{th:ito1}, po lokalizacji na zwartych obszarach.
Dla wszystkich $i,j$ mamy z polaryzacji wariacji
kwadratowej
\[
[X^i,X^j]_t
=\sum_{k=1}^m\int_0^t\sigma_s^{ik}\sigma_s^{jk}\,\dd s.
\]
Istotnie, niezależność współrzędnych $W$ daje
$[W^k,W^\ell]_t=\delta_{k\ell}t$, a izometria Itô,
lokalizacja i aproksymacja współczynników procesami
prostymi pozwalają wyprowadzić tę samą relację dla
całek stochastycznych. Zastosowanie ważonej wersji
lematu \ref{lem:weighted-qv} do sum
$\sum F_{x_ix_j}(t_k,X_{t_k})\Delta_kX^i\Delta_kX^j$
(dla $i\ne j$ przez polaryzację) daje dokładnie
składnik drugiego rzędu z macierzą $\sigma\sigma^{\mathsf T}$.
Część pierwszego rzędu zbiega do całek względem czasu
i procesów Wienera z izometrii i lokalizacji, jak w
dowodzie jednowymiarowym. Otrzymujemy tezę.
\end{proof}

\begin{przyklad}[Ruch Browna i funkcje harmoniczne]
Niech $B$ będzie standardowym $d$-wymiarowym ruchem
Browna, $d\ge2$, i $h\in C^2(\R^d)$ spełnia
$\Delta h=0$. Stosując \eqref{eq:ito-multi} dla
$X=B$, $b=0$, $\sigma=I$, otrzymujemy
$h(B_t)=h(B_0)+\int_0^t\nabla h(B_s)\cdot\dd B_s$.
Jest to lokalny martyngał; po zatrzymaniu przy wyjściu
z kuli staje się kwadratowo całkowalnym martyngałem,
gdyż $\nabla h$ jest ograniczony na odpowiednim
zwartym obszarze. Ten przykład tłumaczy związek
ruchu Browna z funkcjami harmonicznymi.
\end{przyklad}

\begin{zadanie}[Kwadrat odległości w $\R^d$]
Dla $d$-wymiarowego ruchu Browna $B$ wyprowadzić
równanie dla $R_t^2=\norm{B_t}^2$ i obliczyć jego
wartość oczekiwaną, jeśli $B_0=x$ jest deterministyczny.

\emph{Rozwiązanie.} Z \eqref{eq:ito-multi} dla
$F(x)=\sum_i x_i^2$ dostajemy
\[
\norm{B_t}^2=\norm{x}^2+2\int_0^t B_s\cdot\dd B_s+d\,t.
\]
Po wzięciu wartości oczekiwanych mamy
$\E\norm{B_t}^2=\norm{x}^2+dt$. Ostatnia równość
może też zostać otrzymana przez zsumowanie
$\E(B_t^i)^2=x_i^2+t$ dla poszczególnych współrzędnych.
\end{zadanie}

\section{Czas lokalny procesu Wienera: pierwszy kontakt}
Wzór Itô wymaga klasy $C^{1,2}$, tymczasem
funkcja $x\mapsto|x-a|$ nie jest różniczkowalna w $a$.
W przybliżeniu jej druga pochodna koncentruje się
w punkcie $a$. Wartość graniczna daje nowy obiekt:
\emph{czas lokalny}, który mierzy intensywność
przebywania procesu w otoczeniu danego poziomu.

\begin{twierdzenie}[Wzór Tanaki i czas lokalny w ustalonym poziomie]
\label{th:tanaka}
Niech $W$ będzie standardowym procesem Wienera, $a\in\R$.
Przyjmujemy $\operatorname{sgn}(x)=1$ dla $x>0$,
$-1$ dla $x<0$ oraz $0$ dla $x=0$. Istnieje ciągły,
niemalejący, nieujemny proces $(L_t^a)_{0\le t\le T}$,
zaczynający się w zerze, taki że prawie na pewno,
jednocześnie dla wszystkich $t\in[0,T]$,
\[
|W_t-a|=|W_0-a|+\int_0^t\operatorname{sgn}(W_s-a)
\,\dd W_s+L_t^a.
\tag{3.7}\label{eq:tanaka}
\]
Ponadto dla każdego $t$ zachodzi zbieżność w
prawdopodobieństwie
\[
\frac1{2\varepsilon}\int_0^t
\ind{\{|W_s-a|<\varepsilon\}}\,\dd s
\ \longrightarrow\ L_t^a
\qquad(\varepsilon\downarrow0).
\tag{3.8}\label{eq:local-time-occ}
\]
Proces $L^a$ nazywamy \emph{czasem lokalnym} ruchu Browna
na poziomie $a$, w przyjętej normalizacji.
\end{twierdzenie}
\begin{proof}
Zaczynamy od funkcji wypukłej
\[
f_\varepsilon(x)=
\begin{cases}
|x-a|,&|x-a|\ge\varepsilon,\\
\displaystyle\frac{(x-a)^2}{2\varepsilon}+
\frac\varepsilon2,&|x-a|<\varepsilon.
\end{cases}
\]
Jest ona klasy $C^1$, jej pochodna jest Lipschitzowska,
a druga pochodna istnieje poza dwoma punktami
$a\pm\varepsilon$ i wynosi
$f_\varepsilon''(x)=\varepsilon^{-1}
\ind{\{|x-a|<\varepsilon\}}$.
Aby legalnie stosować wzór Itô, przybliżamy
$f_\varepsilon$ splotami z gładkim jądrem
mollifikującym. Pochodne pierwszego rzędu zbiegają
jednostajnie, a drugie — prawie wszędzie i są
jednostajnie ograniczone (przy ustalonym
$\varepsilon$). Ponieważ rozkład $W_s$ ma gęstość
przy każdym $s>0$, a $\int_0^T\Prob(W_s=a\pm\varepsilon)
\,\dd s=0$, zbieżność zdominowana daje
zbieżność składników drugiego rzędu w $L^1(\Omega)$.
Izometria Itô daje zbieżność całek z pochodnych
pierwszego rzędu w $L^2(\Omega)$.
Przechodząc do granicy w gładkim wzorze Itô, uzyskujemy
\[
f_\varepsilon(W_t)=f_\varepsilon(W_0)
+\int_0^t f_\varepsilon'(W_s)\,\dd W_s
+\frac1{2\varepsilon}\int_0^t
\ind{\{|W_s-a|<\varepsilon\}}\,\dd s.
\tag{3.9}\label{eq:tanaka-eps}
\]
Teraz $|f_\varepsilon(x)-|x-a||\le\varepsilon/2$ dla
każdego $x$, zaś $f_\varepsilon'(x)\to
\operatorname{sgn}(x-a)$ poza punktem $a$, przy czym
$|f_\varepsilon'|\le1$. Ponieważ
$\E\int_0^T\ind{\{W_s=a\}}\,\dd s=0$,
z twierdzenia Fubiniego i zbieżności zdominowanej
mamy
$\E\int_0^T|f_\varepsilon'(W_s)-
\operatorname{sgn}(W_s-a)|^2\,\dd s\to0$.
Izometria i nierówność Dooba zapewniają zbieżność
całek stochastycznych jednostajnie w $L^2$.
Z \eqref{eq:tanaka-eps} wynika więc zbieżność ucp
nieujemnych i niemalejących procesów z lewej strony
\eqref{eq:local-time-occ} do ciągłego procesu
\[
L_t^a:=|W_t-a|-|W_0-a|
-\int_0^t\operatorname{sgn}(W_s-a)\,\dd W_s.
\]
Wybieramy podciąg zbieżny jednostajnie prawie na pewno;
jego granica jest niemalejąca i nieujemna. Dostajemy
\eqref{eq:tanaka} oraz \eqref{eq:local-time-occ}.
\end{proof}

\begin{zadanie}[Oczekiwany czas lokalny w zerze]
Dla procesu Wienera $W_0=0$ obliczyć $\E L_t^0$.

\emph{Rozwiązanie.} Biorąc wartości oczekiwane w
\eqref{eq:tanaka}, zauważamy, że całka
$\int_0^t\operatorname{sgn}(W_s)\,\dd W_s$ jest
kwadratowo całkowalnym martyngałem o średniej zero,
gdyż integrand jest ograniczony. Zatem
$\E L_t^0=\E|W_t|=\sqrt{2t/\pi}$.
Wartość jest dodatnia, mimo że $\Prob(W_s=0)=0$
dla każdej ustalonej chwili $s>0$: czas lokalny
mierzy nagromadzenie wizyt \emph{w dowolnie małych
otoczeniach} zera, a nie miarę Lebesgue'a zbioru
chwil, w których $W_s=0$.
\end{zadanie}

\begin{uwaga}[Czego jeszcze nie dowiedliśmy?]
W tym rozdziale skonstruowaliśmy $L_t^a$ dla każdego
\emph{ustalonego} $a$. Wspólnie ciągłą wersję w
parametrach $(t,a)$, pełny wzór okupacyjny i
analogiczne wyniki dla ogólnych półmartyngałów
przedstawimy po wprowadzeniu odpowiednich narzędzi.
Nie wolno uważać \eqref{eq:local-time-occ} za
jednostajną w $a$ zbieżność bez dodatkowego dowodu.
\end{uwaga}


% ============================================================================
% ROZDZIAL 4 — STRATONOVICH
% ============================================================================
